\documentclass{report}
\usepackage{amsmath,amssymb}
\setlength{\parindent}{0mm}
\setlength{\parskip}{1em}
\begin{document}
\begin{center}
\rule{6in}{1pt} \
{\large Hormozd Gahvari \\
{\bf Embedding a Performance Model into hypre's BoomerAMG Solver}}

Lawrence Livermore National Laboratory \\ P O Box 808 \\ L-561 \\ Livermore \\ CA 94551
\\
{\tt gahvari1@llnl.gov}\\
William Gropp\\
Kirk E. Jordan\\
Martin Schulz\\
Ulrike M. Yang\end{center}

Agglomerating data on coarse grids is a commonly used technique in
multigrid solvers to reduce communication and improve performance on
parallel machines. In our past work, we showed that we can use a
performance model at runtime to guide data agglomeration in algebraic
multigrid to ensure that it is done in a way that improves performance.
However, performance models require measurements of machine parameters
and an idea of the factors that affect performance on each machine, which
makes incorporating a performance model into production software a
difficult task. In this talk, we discuss how we incorporated a
performance model into the BoomerAMG solver in the hypre library to
control data agglomeration on coarse grids. The model effectively makes
decisions on when in the multigrid cycle to agglomerate data, and how
much agglomeration to perform. It bases its decisions on runtime
performance measurements, and conveniently allows for reuse of the
measured parameters and decision information in future solves to save
time. This results in improved performance and scalability while also
leaving the user free of having to deal with all of the underlying
machine details.


\end{document}
